\documentclass[11pt,a4paper]{article}

\usepackage[utf8]{inputenc}
\usepackage[T1]{fontenc}
\IfFileExists{lmodern.sty}{\usepackage{lmodern}}{}
\IfFileExists{french.ldf}{\usepackage[french]{babel}}{}
\usepackage[margin=2.5cm]{geometry}
\usepackage{amsmath,amssymb,bm}
\usepackage{xcolor}
\usepackage{listings}
\usepackage[most]{tcolorbox}
\usepackage{booktabs,tabularx,longtable}
\usepackage{enumitem}
\usepackage[hidelinks]{hyperref}
\setlength{\emergencystretch}{3em}
\let\usorig\_
\renewcommand{\_}{\usorig\allowbreak}   % coupure possible après _ dans les noms de fonctions

% ---------------------------------------------------------------------------
% Mise en forme du code
% ---------------------------------------------------------------------------
\definecolor{codebg}{RGB}{247,247,249}
\definecolor{codeframe}{RGB}{200,205,215}
\definecolor{kw}{RGB}{0,70,160}
\definecolor{cm}{RGB}{110,120,130}
\definecolor{st}{RGB}{160,60,20}
\definecolor{notebg}{RGB}{235,243,252}
\definecolor{noteframe}{RGB}{70,120,190}

\lstdefinestyle{py}{
  language=Python,
  basicstyle=\ttfamily\footnotesize,
  keywordstyle=\color{kw}\bfseries,
  commentstyle=\color{cm}\itshape,
  stringstyle=\color{st},
  showstringspaces=false,
  breaklines=true,
  breakatwhitespace=true,
  columns=fullflexible,
  keepspaces=true,
  tabsize=4,
  upquote=true,
  literate=
    {é}{{\'e}}1 {è}{{\`e}}1 {ê}{{\^e}}1 {ë}{{\"e}}1
    {à}{{\`a}}1 {â}{{\^a}}1 {ç}{{\c{c}}}1 {ù}{{\`u}}1 {û}{{\^u}}1
    {î}{{\^i}}1 {ï}{{\"i}}1 {ô}{{\^o}}1 {œ}{{\oe}}1
    {É}{{\'E}}1 {È}{{\`E}}1 {À}{{\`A}}1 {Ç}{{\c{C}}}1
    {«}{{\guillemotleft}}1 {»}{{\guillemotright}}1
    {→}{{$\rightarrow$}}1 {≤}{{$\le$}}1 {≥}{{$\ge$}}1 {×}{{$\times$}}1
    {↑}{{$\uparrow$}}1 {↓}{{$\downarrow$}}1
}
\lstset{style=py}

% Bloc de code encadré :  \begin{code}[Titre]{fichier.py} ... \end{code}
\tcbuselibrary{listings}
\newtcblisting{code}[2][]{
  listing only, breakable, enhanced,
  colback=codebg, colframe=codeframe, boxrule=0.4pt, arc=2pt,
  left=4pt, right=4pt, top=2pt, bottom=2pt,
  title={\small\ttfamily #2\if\relax\detokenize{#1}\relax\else\ \textrm{--- #1}\fi},
  coltitle=black, colbacktitle=codeframe!50,
  listing options={style=py}
}

% Encadré « Dans le code » : correspondance théorie -> fichiers, fonctions, paramètres
\newtcolorbox{danslecode}{
  breakable, enhanced, colback=notebg, colframe=noteframe, boxrule=0.5pt, arc=2pt,
  left=6pt, right=6pt, top=4pt, bottom=4pt,
  title={\small\bfseries Dans le code}, coltitle=white, colbacktitle=noteframe
}

\newcommand{\f}[1]{\texttt{#1}}        % nom de fichier / fonction / paramètre

\title{Architecture V-LSTM pour la prévision financière}
\author{Alexandre Berlia}
\date{Octobre 2026}

\begin{document}
\maketitle
\tableofcontents
\bigskip

% ===========================================================================
\section{Introduction}
% ===========================================================================
L'objectif du modèle est d'estimer la probabilité qu'un titre financier enregistre une
hausse à un horizon donné, afin de construire une stratégie de trading systématique.
Le modèle étudié est un modèle de \emph{deep learning} appelé V-LSTM. Son architecture
combine un \emph{Variable Selection Network} (VSN) et un \emph{Long Short-Term Memory} (LSTM).
Le VSN apprend quelles variables sont pertinentes à chaque date et leur attribue des
poids. Le LSTM est un réseau récurrent qui traite une séquence temporelle et conserve
les informations utiles issues des périodes précédentes grâce à des portes d'entrée,
d'oubli et de sortie. Les variables explicatives comprennent notamment des signaux
construits à partir d'indicateurs techniques. Ces signaux sont renseignés dans le fichier
\f{signaux.py}. Il s'agit principalement de variables comprises entre 0 et 1, représentant
la présence et la fraîcheur d'un signal d'achat ou de vente. L'horizon de prédiction est
sélectionné par un algorithme d'optimisation. Cet algorithme compare plusieurs horizons
sur un échantillon de validation et retient celui qui produit les meilleures performances
après prise en compte des coûts de transaction.

% ---------------------------------------------------------------------------
\subsection{Organisation du code}
% ---------------------------------------------------------------------------
Le projet est découpé en modules Python, chacun correspondant à une étape du processus.
Le tableau~\ref{tab:fichiers} donne la correspondance ; les sections suivantes renvoient
aux fichiers et fonctions concernés dans des encadrés «~Dans le code~».

\begin{table}[h]
\centering\small
\begin{tabularx}{\linewidth}{@{}l X l@{}}
\toprule
\textbf{Fichier} & \textbf{Rôle} & \textbf{Section} \\
\midrule
\f{signaux.py}      & Liste des règles de franchissement (indicateur, seuil, sens) et calcul de la fraîcheur & \ref{sec:signaux} \\
\f{data\_prep.py}   & Assemblage prix + indicateurs + score, cible, découpage temporel, normalisation & \ref{sec:notations}, \ref{sec:score}, \ref{sec:decoupage} \\
\f{windows.py}      & Construction des fenêtres de $T$ séances et des \emph{batchs} & \ref{sec:decoupage} \\
\f{models.py}       & Réseau : GLU, GRN, VSN, LSTM, couche de sortie (\f{VSNLSTM}) & \ref{sec:vsn}, \ref{sec:lstm} \\
\f{train.py}        & Entraînement (perte, Adam, \emph{early stopping}), chargement, prédiction & \ref{sec:perte}, \ref{sec:epoques} \\
\f{resultats.py}    & Rapport sur la validation : AUC, calibration, poids du VSN, lecture par régime & \ref{sec:pipeline} \\
\f{horizons.py}     & Comparaison des horizons $H$ sur la validation & \ref{sec:horizons} \\
\f{Portefeuille.py} & Choix des seuils, backtest de la stratégie, décisions du jour & \ref{sec:seuils}, \ref{sec:backtest} \\
\f{prevision.py}    & \f{figer\_seance} : insère le cours du moment comme séance du jour & \ref{sec:backtest} \\
\f{decision.py}     & Fonctions communes (OHLC, ATR, référence, probabilités du jour) & --- \\
\f{script\_xlk.py}  & Script principal : enchaîne toutes les étapes (point d'entrée) & \ref{sec:pipeline} \\
\bottomrule
\end{tabularx}
\caption{Fichiers du projet.}
\label{tab:fichiers}
\end{table}

\noindent Le script principal est le seul fichier à modifier en usage courant. Les
paramètres y sont regroupés dans un bloc \f{RÉGLAGES} en tête de fichier et dans le
dictionnaire \f{reglages}. La section~\ref{sec:pipeline} décrit le script complet et la
section~\ref{sec:parametres} récapitule tous les paramètres modifiables.

% ===========================================================================
\section{Architecture}
% ===========================================================================
\subsection{Paramètres et notations}\label{sec:notations}
On utilise les notations suivantes :
\begin{itemize}[label=---]
  \item $P_{i,t}$ : prix de clôture du titre $i$ à la date $t$.
  \item $H$ : horizon de prédiction, exprimé en nombre de séances boursières.
  \item $y_{i,t}$ : variable cible binaire définie par
  \[
    y_{i,t} = \mathbf{1}_{\{P_{i,t+H} > P_{i,t}\}} =
    \begin{cases} 1, & \text{si } P_{i,t+H} > P_{i,t},\\ 0, & \text{sinon.} \end{cases}
  \]
  \item $\bm{x}_{i,t}$ : vecteur des variables explicatives du titre $i$ à la date $t$. Il
  contient les signaux techniques et les variables de prix ; le score macroéconomique
  entre séparément, comme contexte (section~\ref{sec:score}).
  \item $K$ : nombre de variables explicatives.
  \item $W_h$ : matrice des poids récurrents appliquée à l'état caché.
  \item $\theta$ : ensemble des paramètres entraînables du modèle.
  \item $\hat p_{i,t}$ : probabilité de hausse prédite par le modèle :
  \[
    \hat p_{i,t} = \mathbb{P}\left(P_{i,t+H} > P_{i,t} \mid \bm{x}_{i,t-T+1},\dots,\bm{x}_{i,t}\right).
  \]
  \item $T$ : longueur de la fenêtre temporelle utilisée par le modèle. Par exemple, $T=20$
  signifie que le modèle observe les 20 dernières séances.
\end{itemize}

\begin{danslecode}
La cible est construite dans \f{data\_prep.py} (fonction \f{price\_features}), appelée par
\f{prepare\_dataset}. Les deux paramètres qui la définissent sont passés dans le script :
\begin{itemize}[nosep]
  \item \f{target="direction"} : cible binaire $y_{i,t}$ (avec \f{"return"}, la cible serait le rendement) ;
  \item \f{horizon=HORIZON} : la valeur de $H$, fixée en tête de script.
\end{itemize}
La fenêtre $T$ est le paramètre \f{FENETRE} (\f{window} dans \f{make\_loaders}).
\end{danslecode}

\begin{code}[cible binaire]{data\_prep.py --- price\_features}
log_close = np.log(prices[close_col])                 # log des clôtures
lc  = log_close.groupby(level="ticker", group_keys=False)
fwd = lc.shift(-horizon) - log_close                  # log(P_{t+H} / P_t), par titre
feats["target"] = (fwd > 0).astype(float).where(fwd.notna()) if target == "direction" else fwd
# les H dernières séances n'ont pas encore de cible : elles sont étiquetées "live"
\end{code}

% ---------------------------------------------------------------------------
\subsection{Fonction de perte}\label{sec:perte}
Le modèle réalise une classification binaire. La fonction de perte utilisée est
l'entropie croisée binaire :
\[
  \mathcal{L}(\theta) = -\frac{1}{N}\sum_{k=1}^{N}\left[y_k\ln(\hat p_k) + (1-y_k)\ln(1-\hat p_k)\right],
\]
où $N$ représente le nombre d'observations, $y_k$ la valeur cible et $\hat p_k$ la
probabilité prédite. Une perte faible signifie que le modèle attribue une probabilité
élevée aux événements qui se réalisent effectivement et une probabilité faible aux
événements qui ne se réalisent pas. En pratique, une version numériquement stable de
cette fonction est appliquée directement aux \emph{logits} du modèle.

\begin{danslecode}
\f{TrainConfig(task="proba")} dans le script sélectionne cette perte. Dans \f{train.py},
elle correspond à \f{nn.BCEWithLogitsLoss()} : le réseau sort un logit $z$ et la sigmoïde
est intégrée à la perte (version stable). La probabilité $\hat p = \sigma(z)$ est calculée
par \f{predict}, qui renvoie une colonne \f{proba}.
\end{danslecode}

% ---------------------------------------------------------------------------
\subsection{Variables techniques}\label{sec:signaux}
Les variables techniques représentent principalement des règles de lecture des
indicateurs plutôt que leur niveau brut. Par exemple, le niveau exact du RSI peut être
moins important que le franchissement d'un seuil. Un signal est déclenché le jour où
l'indicateur passe au-dessus (ou en dessous) d'une valeur déterminée :
\[
  \text{franchissement à la hausse en } t \iff I_{t-1} < s \le I_t ,
\]
et symétriquement à la baisse. Le franchissement est ensuite transformé en une variable
de fraîcheur. On définit :
\[
  f_t = \begin{cases} 2^{-d_t/\lambda}, & \text{si } 0 \le d_t \le 10,\\ 0, & \text{si } d_t > 10, \end{cases}
  \qquad \lambda = 2 \text{ (demi-vie)},
\]
où $d_t$ représente le nombre de séances écoulées depuis le dernier déclenchement du
signal. La valeur est divisée par 2 toutes les $\lambda$ séances. Par exemple :
\begin{center}
\begin{tabular}{c|ccccccc}
$d_t$ & 0 & 1 & 2 & 3 & 4 & 5 & $>10$ \\ \hline
$f_t$ & 1 & 0{,}71 & 0{,}5 & 0{,}35 & 0{,}25 & 0{,}18 & 0
\end{tabular}
\end{center}
Le modèle sait ainsi qu'un signal a été déclenché, mais également depuis combien de temps.
Un signal récent reçoit une valeur élevée, tandis qu'un signal ancien voit
progressivement son influence diminuer.

Chaque signal (d'achat comme de vente) est une variable distincte, à valeurs dans $[0,1]$.
Un signal de vente n'est pas codé $-1$ : le réseau apprend lui-même le signe de l'effet de
chaque variable. Garder deux colonnes séparées permet aussi d'accorder une importance
différente à l'entrée et à la sortie d'un même indicateur.

\begin{danslecode}
\textbf{Règles} : liste \f{SIGNAUX} dans \f{signaux.py}. C'est le seul endroit où l'on
ajoute, retire ou modifie une règle. Chaque entrée donne l'identifiant, la colonne de
l'indicateur (\f{col}, qui doit porter exactement le même nom que la clé du dictionnaire
d'indicateurs), le seuil, le sens du franchissement (\f{"up"}/\f{"down"}) et, pour
information, le sens attendu de l'effet (\f{+1}/\f{-1}, non utilisé par le modèle).

\textbf{Fraîcheur} : \f{memoire\_signaux}. Ses deux paramètres se règlent dans
\f{reglages} (ils sont transmis à \f{prepare\_dataset}) :
\f{demi\_vie} ($\lambda$, défaut 2) et \f{max\_jours} (défaut 10).

\textbf{Indicateurs bruts} : avec \f{garder\_niveaux=False} (défaut), le modèle ne voit
que les signaux, pas le niveau du RSI. \f{garder\_niveaux=True} ajoute les niveaux.

\textbf{Variables de prix} : par défaut, \f{prepare\_dataset} ajoute aussi \f{ret\_1d},
\f{ret\_5d}, \f{ret\_20d}, \f{vol\_20d} et \f{volume\_rel}. \f{variables\_prix=False} les retire.
\end{danslecode}

\begin{code}[liste des règles (extrait)]{signaux.py}
SIGNAUX = [
  {"id": "rsi_up45", "nom": "RSI ↑45",  "col": "rsi",    "seuil": 45,   "cross": "up",   "sens": +1},
  {"id": "rsi_dn65", "nom": "RSI ↓65",  "col": "rsi",    "seuil": 65,   "cross": "down", "sens": -1},
  {"id": "bb_up0",   "nom": "BB% ↑0",   "col": "bb_pct", "seuil": 0.0,  "cross": "up",   "sens": +1},
  {"id": "bb_dn095", "nom": "BB% ↓0,95","col": "bb_pct", "seuil": 0.95, "cross": "down", "sens": -1},
  # ... rev_5, stoch, hist (MACD), er : même structure
]

def franchissement(s, seuil, sens):
    veille = s.groupby(level="ticker").shift(1)
    if sens == "up":
        return (veille < seuil) & (s >= seuil)
    return (veille > seuil) & (s <= seuil)

def memoire_signaux(evenements, demi_vie=2.0, max_jours=10):
    ...
    jours = pos - last                       # séances depuis le dernier franchissement
    val = np.power(0.5, jours / demi_vie)    # f = 2^(-d / demi_vie)
    out[c] = val.where(jours <= max_jours, 0.0).fillna(0.0)
\end{code}

\noindent Les noms de colonnes de \f{SIGNAUX} doivent correspondre aux clés du
dictionnaire d'indicateurs construit dans le script (\f{"rsi"}, \f{"bb\_pct"},
\f{"rev\_5"}, \f{"stoch"}, \f{"er"}, \f{"hist"}). Si un nom ne correspond pas, le signal
n'est jamais déclenché et \f{prepare\_dataset} affiche un avertissement.

% ---------------------------------------------------------------------------
\subsection{Variable Selection Network}\label{sec:vsn}
Le VSN apprend l'importance relative de chaque variable à chaque date. Pour chaque
variable $j\in\{1,\dots,K\}$, le réseau produit un score $a^{(j)}_{i,t}$. Une fonction
\emph{softmax} transforme ces scores en poids normalisés :
\[
  \alpha^{(j)}_{i,t} = \frac{\exp\left(a^{(j)}_{i,t}\right)}{\sum_{\ell=1}^{K}\exp\left(a^{(\ell)}_{i,t}\right)}.
\]
Les poids vérifient :
\[
  \alpha^{(j)}_{i,t} \ge 0 \qquad\text{et}\qquad \sum_{j=1}^{K}\alpha^{(j)}_{i,t} = 1.
\]
La représentation transmise au LSTM est une combinaison pondérée des représentations des
variables :
\[
  \tilde{\bm{x}}_{i,t} = \sum_{j=1}^{K}\alpha^{(j)}_{i,t}\,\bm{e}^{(j)}_{i,t},
\]
où $\bm{e}^{(j)}_{i,t}$ représente la transformation de la variable $j$. Les poids
$\alpha^{(j)}_{i,t}$ permettent d'interpréter les variables auxquelles le modèle accorde le
plus d'importance à une date donnée.

\paragraph{Gated Residual Network (GRN).} Les scores $a^{(j)}$ et les transformations
$\bm{e}^{(j)}$ sont produits par des GRN. Pour une entrée $\bm{a}$ et un contexte $\bm{c}$ :
\[
  \bm\eta_2 = \mathrm{ELU}(W_2\bm{a} + W_3\bm{c} + \bm b_2),\quad
  \bm\eta_1 = W_1\bm\eta_2 + \bm b_1,\quad
  \mathrm{GRN}(\bm a,\bm c) = \mathrm{LayerNorm}\big(\bm a + \mathrm{GLU}(\bm\eta_1)\big),
\]
avec $\mathrm{GLU}(\bm\eta) = \sigma(W_4\bm\eta + \bm b_4)\odot(W_5\bm\eta + \bm b_5)$. La
porte du GLU peut s'approcher de 0 : le bloc se réduit alors à l'identité, ce qui
permet au réseau de «~sauter~» une transformation inutile.

\begin{danslecode}
\f{models.py}, classes \f{GLU}, \f{GRN} et \f{VSN}. Dans \f{VSN.forward} :
\f{proj} donne une représentation de dimension \f{d\_model} à chaque variable,
\f{grn\_weights} produit les scores $a^{(j)}$ (puis \f{softmax}), \f{grn\_vars} produit les
$\bm e^{(j)}$ (un GRN par variable). La taille \f{d\_model} se règle dans
\f{TrainConfig(d\_model=16)}. Les poids $\alpha$ sont renvoyés par \f{predict} dans les
colonnes \f{w\_*} et résumés par \f{rapport} (feuille d'importance des variables).
\end{danslecode}

\begin{code}[sélection des variables]{models.py --- VSN.forward}
xi = [proj(x[..., j:j + 1]) for j, proj in enumerate(self.proj)]   # 1 vecteur par variable
flat = torch.cat(xi, dim=-1)
weights = torch.softmax(self.grn_weights(flat, c), dim=-1)         # alpha_j, somme = 1
xi_tilde = torch.stack([grn(e, c_vars) for grn, e in zip(self.grn_vars, xi)], dim=-2)  # e_j
out = (weights.unsqueeze(-1) * xi_tilde).sum(dim=-2)               # x_tilde = somme alpha_j e_j
\end{code}

% ---------------------------------------------------------------------------
\subsection{Score macroéconomique : variable de contexte}\label{sec:score}
Le score macroéconomique décrit le régime de marché dans lequel la prévision est réalisée.
Le score choisi pour ce modèle est le score niveau, et les variables utilisées pour le
calculer sont les suivantes :
\begin{itemize}[nosep,label=-]
  \item Philadelphia Fed Business Outlook ;
  \item Chicago Fed National Activity Index ;
  \item permis de construire américains ;
  \item ISM Manufacturing PMI ;
  \item variation des emplois non agricoles ;
  \item variation de l'indice des prix PCE ;
  \item Chicago PMI ;
  \item taux de chômage américain ;
  \item production industrielle, en variation annuelle.
\end{itemize}
Ces indicateurs sont des indicateurs avancés du PIB américain (consulter le \f{readme
score.md} pour plus d'informations sur les indicateurs avancés et la construction du
score). Le calcul du score est dans \f{cycle-score.py} du dossier \f{update-score}. Tout
ce qui y est renseigné est «~frais~» (et peut ne pas coller avec les exemples et
commentaires du readme). Pour mettre à jour le score, remplacer le fichier
\f{Data-in-value.csv} par les nouvelles données et lancer \f{calcul-score-macro.py}. Le
nouveau fichier est \f{score\_updated.csv}.

Le score fournit au modèle un contexte permettant de distinguer, par exemple, un signal
technique observé dans un régime favorable du même signal observé dans un régime
défavorable : c'est une variable de contexte dans le modèle. Il n'est pas sélectionné
par le VSN comme les autres variables ; il modifie la façon dont le VSN pondère et lit
les signaux.

\paragraph{Délai de publication.} Le score du mois $M$ n'est connu qu'après la publication
de ses composantes. Il n'est utilisé qu'à partir de la fin du mois $M$ plus 30 jours ;
chaque séance reçoit le dernier score déjà publié à cette date.

\begin{danslecode}
\textbf{Chargement} : \f{load\_macro("score\_updated.csv", date\_col="periode\_M")} dans
\f{data\_prep.py}. Le délai se règle avec \f{publication\_lag\_days} (défaut 30). La
fonction ajoute aussi la variation mensuelle \f{score\_niveau\_chg}.

\textbf{Contexte} : \f{regime\_col="score\_niveau"} dans \f{reglages}. Toutes les colonnes
du score passent alors en contexte, et deux indicatrices sont ajoutées :
\f{regime\_bas} (score $<25$) et \f{regime\_eleve} (score $>75$), seuils réglables par
\f{seuils\_regime=(25, 75)}. Sans \f{regime\_col}, le score serait une variable ordinaire
lue par le VSN.

\textbf{Filtre optionnel} : \f{score\_min=25} (et/ou \f{score\_max}) dans \f{reglages}. Les
jours hors domaine sont étiquetés \f{"hors\_score"} : le modèle n'y apprend pas et n'y
prédit pas, donc la stratégie n'y achète pas. \f{filtre\_sur=("val","test","live")}
garde l'apprentissage sur toutes les périodes et n'applique le filtre qu'aux décisions.
Désactivé dans la version actuelle.
\end{danslecode}

\begin{code}[le score comme contexte]{models.py --- VSNLSTM.forward}
c_r = self.grn_regime(self.regime_proj(regime))   # score -> vecteur de contexte c_t
xi, weights = self.vsn(x, c_r)                    # le contexte entre dans le VSN (poids + GRN_j)
...
z = self.grn_out(h[:, -1], c_r[:, -1])            # et dans le GRN de sortie
\end{code}

% ---------------------------------------------------------------------------
\subsection{Découpage temporel}\label{sec:decoupage}
\subsubsection{Les trois ensembles}
Les données sont séparées chronologiquement en trois ensembles :
\begin{enumerate}[nosep]
  \item un ensemble d'entraînement utilisé pour estimer les paramètres ;
  \item un ensemble de validation utilisé pour choisir les hyperparamètres et les seuils de décision ;
  \item un ensemble de test utilisé pour évaluer la performance finale.
\end{enumerate}
Le découpage doit respecter l'ordre temporel :
\[
  \text{Entraînement} \longrightarrow \text{Validation} \longrightarrow \text{Test}.
\]
La cible d'une date $t$ utilise le prix en $t+H$. Les $H$ dernières séances de
l'entraînement (et de la validation) ont donc une cible qui déborde sur la période
suivante : elles sont retirées (\emph{purge}) pour éviter toute fuite d'information.
La normalisation des variables de prix (z-score, écrêté à $\pm 5$) est estimée sur
l'entraînement seulement, puis appliquée telle quelle aux autres périodes.

Un échantillon est une fenêtre des $T=20$ dernières séances d'un seul titre, avec la
fraîcheur de chaque signal et le contexte du score à chaque séance. Le modèle est commun
à toutes les actions : les fenêtres de tous les titres sont mélangées, ce qui multiplie
le nombre d'exemples de chaque signal. Les fenêtres sont présentées au modèle par paquets
de 256, appelés \emph{batchs}.

\begin{danslecode}
\textbf{Dates} : dans \f{reglages} du script.
\begin{itemize}[nosep]
  \item \f{train\_start} : début de l'entraînement (les dates antérieures servent seulement
  à remplir l'historique des premières fenêtres) ;
  \item \f{train\_end} : fin de l'entraînement ; \f{val\_end} : fin de la validation ;
  \item tout ce qui suit \f{val\_end} est le test ; les $H$ dernières séances, sans cible,
  sont étiquetées \f{"live"} et servent aux décisions du jour.
\end{itemize}
Le découpage est fait par \f{time\_split} (\f{data\_prep.py}) ; l'étiquette de chaque ligne
est dans la colonne \f{split} de \f{data.df}.

\textbf{Fenêtres et batchs} : \f{make\_loaders(data, window=FENETRE, batch\_size=256)}
dans \f{windows.py}. Une fenêtre est rejetée si elle contient un trou de plus de 7 jours
calendaires (\f{max\_gap\_days}).
\end{danslecode}

\begin{code}[préparation et fenêtres]{script\_xlk.py}
reglages = dict(train_start="2005-01-01", train_end="2016-12-31", val_end="2019-12-31",
                horizon=HORIZON, target="direction", signaux=True,
                indicator_field_level=0, regime_col="score_niveau")

data    = prepare_dataset(prices, indicateurs, score, **reglages)
loaders = make_loaders(data, window=FENETRE, batch_size=256)   # {"train", "val", "test"}
print(data.df["split"].value_counts())                         # contrôle du découpage
\end{code}

\subsubsection{Les époques, la validation et l'arrêt}\label{sec:epoques}
Une époque est un passage complet sur toutes les fenêtres d'entraînement. Avec 30 titres
sur 12 ans, on obtient environ 90\,000 fenêtres, soit environ 350 batchs, c'est-à-dire
350 ajustements des paramètres par époque.

Entraînement et validation alternent : à la fin de chaque époque, on mesure l'erreur de
prédiction sur la période de validation, sans modifier les paramètres. Si l'erreur est la
plus basse observée, les paramètres sont mis en réserve. L'erreur de validation est
d'abord décroissante, puis croissante quand le modèle commence à surapprendre. Après 8
époques consécutives sans amélioration, on arrête et on retient les paramètres mis en
réserve. L'optimiseur est Adam ; le pas d'apprentissage est divisé par 2 après 4 époques
sans amélioration, et la norme du gradient est plafonnée à 1 (\emph{gradient clipping}).

\begin{danslecode}
\f{TrainConfig} dans \f{train.py}, instancié dans le script :
\begin{itemize}[nosep]
  \item \f{patience=8} : époques sans amélioration avant l'arrêt ;
  \item \f{max\_epochs=100} : plafond d'époques ;
  \item \f{lr=1e-3}, \f{lr\_patience=4}, \f{lr\_factor=0.5} : pas d'apprentissage et sa réduction ;
  \item \f{dropout=0.2}, \f{weight\_decay=1e-5} : régularisation ;
  \item \f{grad\_clip=1.0} ; \f{seed=42} (reproductibilité).
\end{itemize}
\f{train(loaders, cfg, save\_path="modele\_proba.pt")} enregistre les meilleurs poids
dans le fichier \f{.pt}. \f{load\_model("modele\_proba.pt")} les recharge sans réentraîner.
\end{danslecode}

\begin{code}[entraînement]{script\_xlk.py}
cfg = TrainConfig(task="proba", d_model=16, dropout=0.2, patience=8)
model, history = train(loaders, cfg, save_path=FICHIER_MODELE)   # .pt = meilleurs poids
rapport(FICHIER_MODELE, loaders, split="val", out_dir="resultats_proba", data=data)
\end{code}

% ---------------------------------------------------------------------------
\subsection{Choix de l'horizon}\label{sec:horizons}
L'horizon $H$ est choisi en comparant plusieurs valeurs sur la validation : pour chaque
horizon, on prépare les données, on entraîne un modèle (régression logistique rapide ou
réseau complet), puis on évalue l'AUC et la stratégie après frais.

\begin{danslecode}
\f{comparer\_horizons} dans \f{horizons.py}. Les horizons testés se passent dans
\f{horizons=(...)} ; \f{methode="logistique"} (rapide) ou \f{"reseau"} (VSN-LSTM, plusieurs
minutes par horizon). Les autres arguments de \f{prepare\_dataset} sont transmis tels quels
(\f{**reglages}, sans \f{horizon}). L'horizon retenu est ensuite reporté dans
\f{HORIZON} en tête du script. Attention : cette fonction évalue la stratégie avec le
backtest simple de \f{decision.py} (détention fixe de $H$ séances), pas avec la stratégie
portefeuille de la section~\ref{sec:backtest}.
\end{danslecode}

\begin{code}[comparaison des horizons]{horizons.py}
from horizons import comparer_horizons
r = {k: v for k, v in reglages.items() if k != "horizon"}
tab = comparer_horizons(prices, indicateurs, score, horizons=(5, 10, 15, 22, 30),
                        methode="logistique", window=FENETRE, couts_bps=5,
                        entree="cloture", **r)
print(tab.to_string())
\end{code}

% ---------------------------------------------------------------------------
\subsection{Seuils et backtest}
\subsubsection{Choix des seuils}\label{sec:seuils}
Le modèle figé calcule ses probabilités sur la période de validation. La fonction
\f{choisir\_seuils\_portefeuille} essaie plusieurs couples de seuils candidats (achat,
sortie) ; pour chacun, elle simule la stratégie complète et classe les couples par
ratio de Sharpe net de frais (rendement rapporté à la volatilité). Les seuils d'achat
candidats sont, par défaut, les quantiles 70\,\%, 80\,\%, 90\,\% et 95\,\% des probabilités
de validation ; les seuils de sortie sont 0,50, la médiane des probabilités et
«~aucun~» (sortie par le stop seulement). Le script retient la première ligne du
tableau ; il faut vérifier qu'elle se trouve sur un plateau (des seuils voisins donnent
des résultats proches) et non sur un pic isolé. Le test n'est regardé qu'une fois, avec
les seuils figés.

\begin{danslecode}
Fonction \f{choisir\_seuils\_portefeuille} (\f{Portefeuille.py}). Paramètres :
\f{grille\_achat}, \f{grille\_sortie} (listes de seuils), \f{min\_trades} (nombre minimal de
trades pour qu'une combinaison soit «~évaluable~», 10 dans le script), plus les mêmes
options que le backtest (\f{**OPTIONS}). Les seuils retenus sont enregistrés dans
\f{seuils.json} et relus pour les décisions du jour. Une grille grossière limite le
surapprentissage : plus on teste de combinaisons, plus la meilleure est flatteuse.
\end{danslecode}

\begin{code}[choix des seuils sur la validation]{script\_xlk.py}
grille = choisir_seuils_portefeuille(predict(model, loaders["val"]), prices,
                                     min_trades=10, **OPTIONS)
# grille personnalisée : ajouter grille_achat=[0.55, 0.6, 0.65], grille_sortie=[0.4, 0.5]
print(grille.head(10).round(3).to_string())
seuil_achat  = float(grille.iloc[0]["seuil_achat"])
ss           = grille.iloc[0]["seuil_sortie"]
seuil_sortie = None if isinstance(ss, str) else float(ss)        # "aucun (stop seul)"
json.dump({"seuil_achat": seuil_achat, "seuil_sortie": seuil_sortie}, open(FICHIER_SEUILS, "w"))
\end{code}

\subsubsection{Règles de la stratégie et backtest}\label{sec:backtest}
La prise de décision est la suivante. Lorsqu'on lance la simulation, on compare
$\hat p_{i,t}$ au seuil d'achat. Si la probabilité est supérieure, on prend la position
au dernier cours disponible : si l'on simule le modèle à 17\,h, le cours qui sert à
calculer la probabilité et à acheter est le même, celui de 17\,h. Cela évite le
\emph{gap} d'ouverture du lendemain. Dans le backtest, ce cours est la clôture.

Le portefeuille peut détenir au maximum deux fois plus de positions que de titres (et au
plus 10), avec au plus deux positions par titre dont les signaux d'achat sont séparés
d'au moins 3 séances (exemple : aujourd'hui j'ai un signal sur Schneider ; 4 séances plus
tard j'en ai un autre, sans signal de vente entre-temps : je rachète une deuxième ligne).
Au plus deux achats sont faits par jour, en priorité les plus fortes probabilités.

Un titre sort :
\begin{itemize}[nosep]
  \item lorsque sa probabilité de hausse est inférieure au seuil de sortie : toutes ses
  lignes sont vendues ;
  \item ou, si le stop est activé, lorsque le cours touche $\text{prix d'achat} - k\times\text{ATR}$
  (calculé ligne par ligne, avec l'ATR du jour d'achat).
\end{itemize}
Il n'y a pas de sortie au bout d'une durée fixe. Le capital est découpé en autant de
tranches que \f{max\_positions} : chaque achat investit $1/\f{max\_positions}$ de la valeur
courante du portefeuille, le reste attend en liquidités (rémunérées à 0\,\%). Les frais
(\f{couts\_bps}) sont prélevés à l'achat et à la vente. La référence est l'achat et la
conservation des mêmes titres, équipondérés, frais compris.

\begin{danslecode}
\textbf{Options de la stratégie} : dictionnaire renvoyé par la fonction \f{options} du
script, transmis tel quel au choix des seuils, au backtest et aux décisions du jour.
\begin{itemize}[nosep]
  \item \f{max\_positions} : nombre de lignes simultanées (taille d'une ligne $=1/$\f{max\_positions}) ;
  \item \f{max\_par\_titre=2}, \f{ecart\_min=3} : lignes par titre et écart minimal entre deux achats ;
  \item \f{max\_entrees\_jour=2} : achats par jour ;
  \item \f{stop\_atr=2}, \f{atr=...} : stop à 2 ATR ; \f{stop\_atr=None} le désactive ;
  \item \f{duree\_max=None} : pas de sortie au temps (un entier la réactive) ;
  \item \f{couts\_bps=5} : frais par transaction, en points de base ;
  \item \f{entree="cloture"} : exécution au cours qui a servi au calcul (\f{"ouverture"} :
  à l'ouverture suivante) ;
  \item \f{recroisement=True} (optionnel) : la 2\textsuperscript{e} ligne exige que la
  probabilité soit repassée sous le seuil d'achat entre-temps.
\end{itemize}
\textbf{Backtest} : \f{backtest\_portefeuille} puis \f{rapport\_portefeuille}
(Excel + graphique dans \f{out\_dir}).
\end{danslecode}

\begin{code}[options et backtest sur le test]{script\_xlk.py}
def options(atr):
    return dict(couts_bps=5, entree="cloture",
                stop_atr=2, atr=atr,                     # stop_atr=None : stop désactivé
                max_par_titre=2, ecart_min=3,
                max_positions=min(10, 2 * len(prices)),  # chaque ligne = 1/max_positions
                max_entrees_jour=2,
                duree_max=None)                          # sortie : signal ou stop

bt = backtest_portefeuille(predict(model, loaders["test"]), prices,
                           seuil_achat, seuil_sortie, **OPTIONS)
print(rapport_portefeuille(bt, out_dir="backtest_etf_tech").round(3).to_string())
\end{code}

\paragraph{Lecture du backtest.} Le rendement annuel de la stratégie doit être comparé à
celui de la référence en tenant compte de l'exposition : une stratégie investie 30\,\% du
temps sur un actif qui fait 19\,\% par an rapporte environ 5,7\,\% si elle choisit ses jours
au hasard. Les critères pertinents sont le ratio de Sharpe net, la perte maximale, la
part du temps investie (\f{exposition moyenne}) et le rendement rapporté à cette
exposition.

% ===========================================================================
\section{Le modèle récurrent}\label{sec:lstm}
% ===========================================================================
\subsection{RNN et disparition du gradient}
Un réseau de neurones récurrent, ou RNN, lit une séquence pas de temps par pas de temps.
À chaque date, il met à jour un état caché $\bm h_t$ qui résume les informations
observées dans le passé.

\subsubsection{Équation du RNN}
L'état caché est défini par :
\[
  \bm h_t = \tanh\left(W_h\bm h_{t-1} + W_x\bm x_t + \bm b_h\right),
\]
où :
\begin{itemize}[nosep,label=---]
  \item $\bm x_t$ est le vecteur des variables à la date $t$ ;
  \item $\bm h_{t-1}$ est l'état caché de la date précédente ;
  \item $W_x$ est la matrice des poids appliquée aux variables ;
  \item $W_h$ est la matrice des poids récurrents ;
  \item $\bm b_h$ est le vecteur de biais ;
  \item $\tanh$ est la fonction d'activation tangente hyperbolique.
\end{itemize}
La prédiction dépend donc à la fois de l'information actuelle $\bm x_t$ et du résumé du
passé $\bm h_{t-1}$.

\subsubsection{Rétropropagation dans le temps}
Pour entraîner un RNN, le réseau est déroulé dans le temps. Chaque pas de temps peut être
interprété comme une couche d'un réseau profond. Le gradient de la perte finale
$\mathcal L_T$ par rapport à un état caché lointain $\bm h_k$ s'écrit :
\[
  \frac{\partial\mathcal L_T}{\partial\bm h_k} = \frac{\partial\mathcal L_T}{\partial\bm h_T}
  \prod_{t=k+1}^{T}\frac{\partial\bm h_t}{\partial\bm h_{t-1}}.
\]
Pour un RNN utilisant la fonction $\tanh$, le jacobien local est :
\[
  \frac{\partial\bm h_t}{\partial\bm h_{t-1}} = \mathrm{diag}\left(1-\bm h_t^2\right)W_h.
\]
La rétropropagation doit donc multiplier un grand nombre de matrices lorsque la distance
$T-k$ est importante.

\subsubsection{Disparition du gradient}
Supposons que la norme de chaque jacobien soit majorée par une constante $\gamma$ :
\[
  \left\|\frac{\partial\bm h_t}{\partial\bm h_{t-1}}\right\| \le \gamma.
\]
On obtient alors :
\[
  \left\|\frac{\partial\mathcal L_T}{\partial\bm h_k}\right\| \le
  \left\|\frac{\partial\mathcal L_T}{\partial\bm h_T}\right\|\gamma^{T-k}.
\]
Si $0<\gamma<1$, alors $\gamma^{T-k}$ converge exponentiellement vers zéro lorsque la
distance temporelle $T-k$ augmente. Le gradient reçu par les états anciens devient donc
presque nul. Le réseau ne parvient plus à modifier correctement les poids responsables des
informations lointaines. Ce phénomène est appelé \textbf{disparition du gradient}. À
l'inverse, si $\gamma>1$, le produit peut devenir très grand. On parle alors
d'\textbf{explosion du gradient}, que l'on limite en plafonnant la norme du gradient
(\f{grad\_clip} dans \f{TrainConfig}).

\subsection{Long Short-Term Memory}
Le LSTM est une extension du RNN conçue pour mieux conserver les dépendances temporelles
longues. En plus de l'état caché $\bm h_t$, le LSTM introduit un état de cellule $\bm c_t$.
Cet état constitue la mémoire interne de long terme du réseau. La mémoire est mise à jour
par une combinaison d'additions et de produits élément par élément. Cette structure crée
un chemin plus direct pour la propagation du gradient. À chaque date, le LSTM calcule
successivement :
\begin{enumerate}[nosep]
  \item une porte d'oubli ;
  \item une porte d'entrée ;
  \item un candidat mémoire ;
  \item une porte de sortie.
\end{enumerate}
On définit le vecteur concaténé :
\[
  \bm z_t = \begin{bmatrix}\bm h_{t-1}\\ \bm x_t\end{bmatrix}.
\]

\subsubsection{Porte d'oubli}
La porte d'oubli est définie par :
\[
  \bm f_t = \sigma\left(W_f\bm z_t + \bm b_f\right),
\]
où $\sigma$ est la fonction sigmoïde :
\[
  \sigma(x) = \frac{1}{1+e^{-x}}.
\]
Chaque composante de $\bm f_t$ appartient à l'intervalle $[0,1]$.
\begin{itemize}[nosep,label=---]
  \item Si une composante est proche de 1, l'information correspondante est conservée.
  \item Si une composante est proche de 0, l'information correspondante est effacée.
\end{itemize}
Dans un contexte financier, cette porte peut permettre au modèle d'oublier un régime de
marché devenu obsolète après une rupture macroéconomique ou un changement important de
volatilité.

\subsubsection{Porte d'entrée}
La porte d'entrée détermine la quantité de nouvelle information qui doit être ajoutée à
la mémoire :
\[
  \bm i_t = \sigma\left(W_i\bm z_t + \bm b_i\right).
\]
Une valeur proche de 1 autorise l'enregistrement de la nouvelle information. Une valeur
proche de 0 empêche cette information de modifier la mémoire. Par exemple, le modèle peut
décider d'accorder une importance particulière à un signal apparu le jour de la
publication des résultats d'une entreprise.

\subsubsection{Candidat mémoire}
Le candidat mémoire représente la nouvelle information susceptible d'être enregistrée :
\[
  \tilde{\bm c}_t = \tanh\left(W_c\bm z_t + \bm b_c\right).
\]
Comme la fonction $\tanh$ prend ses valeurs dans l'intervalle $[-1,1]$, le candidat
mémoire contient à la fois :
\begin{itemize}[nosep,label=---]
  \item le signe de l'information ;
  \item l'intensité de l'information.
\end{itemize}
Dans le cas financier, une valeur positive peut représenter une information favorable,
tandis qu'une valeur négative peut représenter une information défavorable.

\subsubsection{Mise à jour de la cellule mémoire}
La nouvelle mémoire est calculée par :
\[
  \bm c_t = \bm f_t\odot\bm c_{t-1} + \bm i_t\odot\tilde{\bm c}_t,
\]
où $\odot$ représente le produit élément par élément. Cette équation contient deux parties :
\[
  \underbrace{\bm f_t\odot\bm c_{t-1}}_{\text{mémoire ancienne conservée}}
  \;+\;
  \underbrace{\bm i_t\odot\tilde{\bm c}_t}_{\text{nouvelle information ajoutée}}.
\]
Le LSTM peut donc conserver une information ancienne tout en ajoutant sélectivement de
nouvelles informations.

\subsubsection{Porte de sortie}
La porte de sortie détermine la partie de la mémoire qui doit être transmise à l'état
caché :
\[
  \bm o_t = \sigma\left(W_o\bm z_t + \bm b_o\right).
\]
L'état caché est ensuite calculé par :
\[
  \bm h_t = \bm o_t\odot\tanh(\bm c_t).
\]
L'état $\bm h_t$ constitue la représentation visible du LSTM à la date $t$, tandis que
$\bm c_t$ représente sa mémoire interne.

\subsubsection{Résumé des équations du LSTM}
Les équations complètes peuvent être regroupées ainsi :
\begin{align}
  \bm f_t &= \sigma\left(W_f\bm z_t + \bm b_f\right),\\
  \bm i_t &= \sigma\left(W_i\bm z_t + \bm b_i\right),\\
  \tilde{\bm c}_t &= \tanh\left(W_c\bm z_t + \bm b_c\right),\\
  \bm c_t &= \bm f_t\odot\bm c_{t-1} + \bm i_t\odot\tilde{\bm c}_t,\\
  \bm o_t &= \sigma\left(W_o\bm z_t + \bm b_o\right),\\
  \bm h_t &= \bm o_t\odot\tanh(\bm c_t).
\end{align}

\subsubsection{Pourquoi le LSTM limite la disparition du gradient}
La dérivée de la cellule mémoire par rapport à la cellule précédente contient
directement la porte d'oubli :
\[
  \frac{\partial\bm c_t}{\partial\bm c_{t-1}} = \bm f_t.
\]
Sur plusieurs périodes, le gradient contient donc le produit :
\[
  \frac{\partial\bm c_T}{\partial\bm c_k} = \prod_{t=k+1}^{T}\bm f_t.
\]
Lorsque les composantes des portes d'oubli restent proches de 1, le gradient peut
circuler sur une longue période sans disparaître immédiatement. Le LSTM n'élimine pas
totalement le problème de disparition du gradient, mais son architecture lui permet de
mieux préserver les dépendances longues qu'un RNN simple.

\begin{danslecode}
Le LSTM n'est pas recodé : c'est \f{nn.LSTM} de PyTorch, créé dans \f{VSNLSTM.\_\_init\_\_}
(\f{models.py}). Il reçoit la sortie du VSN, de dimension \f{d\_model}. Le nombre de
couches empilées se règle par \f{TrainConfig(n\_layers=1)}. Après le LSTM, un bloc
\f{GateAddNorm} ajoute un saut résiduel : si le LSTM n'apporte rien, la sortie du VSN
passe directement.
\end{danslecode}

\subsection{Couche de prédiction}
Après avoir traité une fenêtre de $T$ séances, le modèle utilise le dernier état caché
$\bm h_T$, enrichi par un dernier GRN qui reçoit le contexte du score, pour produire un
logit :
\[
  z_{i,t} = \bm w_p^\top\,\mathrm{GRN}(\bm h_{i,T},\bm c_T) + b_p.
\]
La probabilité de hausse est obtenue avec une fonction sigmoïde :
\[
  \hat p_{i,t} = \sigma(z_{i,t}) = \frac{1}{1+\exp(-z_{i,t})}.
\]
La décision de trading peut ensuite être définie à partir d'un seuil d'entrée $s_{\text{achat}}$ :
\[
  \text{Décision}_{i,t} = \begin{cases}\text{acheter}, & \text{si } \hat p_{i,t}\ge s_{\text{achat}},\\
  \text{ne pas acheter}, & \text{sinon.}\end{cases}
\]
Le seuil ne doit pas nécessairement être égal à 0,5. Il est sélectionné sur l'échantillon
de validation afin de maximiser une mesure de performance après prise en compte des
coûts de transaction (section~\ref{sec:seuils}).

\subsection{Synthèse de l'architecture V-LSTM}
Le fonctionnement général du modèle peut être résumé par la chaîne suivante :
\[
  \text{Variables brutes}\longrightarrow\text{signaux et fraîcheur}\longrightarrow\text{VSN}
  \longrightarrow\text{LSTM}\longrightarrow\text{couche linéaire}\longrightarrow\hat p_{i,t}.
\]
À chaque séance de la fenêtre :
\begin{itemize}[nosep,label=-]
  \item le score est transformé en un vecteur de contexte $\bm c_t$ ;
  \item le VSN calcule un poids pour chaque signal, $\bm v_t = \mathrm{Softmax}(\mathrm{GRN}(\Xi_t,\bm c_t))$ :
  poids positifs, de somme 1, recalculés chaque jour selon les signaux et le score ;
  \item chaque signal est «~lu~» par son propre GRN, qui reçoit lui aussi le contexte :
  $\tilde\xi^{(j)}_{t} = \mathrm{GRN}_j(\xi^{(j)}_{t},\bm c_t)$ ; c'est ici qu'un franchissement
  peut être interprété différemment selon le score ;
  \item les 20 séances résumées passent dans un LSTM qui capte l'enchaînement des signaux
  dans le temps ;
  \item un dernier GRN, enrichi du contexte, produit un vecteur $\bm z$, puis la probabilité
  $\hat p = \sigma(\bm w^\top\bm z + b)$.
\end{itemize}
Le VSN répond à la question :
\begin{center}\emph{Quelles variables sont importantes à cette date ?}\end{center}
Le LSTM répond à la question :
\begin{center}\emph{Comment ces informations ont-elles évolué au cours des dernières séances ?}\end{center}
La couche finale transforme ensuite la représentation temporelle en une probabilité de
hausse à l'horizon $H$.

\begin{code}[la chaîne complète du réseau]{models.py --- VSNLSTM.forward}
c_r = self.grn_regime(self.regime_proj(regime))   # 1. score -> contexte c_t        (B, T, d)
xi, weights = self.vsn(x, c_r)                    # 2. VSN : poids + GRN_j          (B, T, d)
h, _ = self.lstm(xi, state)                       # 3. LSTM sur les T séances       (B, T, d)
h = self.post_lstm(h, residual=xi)                #    saut résiduel à porte
z = self.grn_out(h[:, -1], c_r[:, -1])            # 4. dernière séance + contexte   (B, d)
pred = self.head(z)                               # 5. logit ; proba = sigmoid(logit)
return {"prediction": pred, "vsn_weights": weights}
\end{code}

% ===========================================================================
\section{Mise en œuvre : le script principal}\label{sec:pipeline}
% ===========================================================================
Le script \f{script\_xlk.py} enchaîne toutes les étapes. Il se compose de trois parties :
les réglages, le bloc 1 (exécuté une fois) et le bloc 2 (exécuté chaque jour).

\subsection{Réglages}
\begin{code}[en-tête du script]{script\_xlk.py}
ENTRAINER = True            # True : bloc 1 + bloc 2 ; False : bloc 2 seulement
TICKERS   = ["XLK"]         # univers de titres
DATE      = "2005-01-01"    # début du téléchargement des prix
HORIZON   = 22              # H, en séances
FENETRE   = 20              # T, en séances
FICHIER_MODELE = "modele_proba.pt"
FICHIER_SEUILS = "seuils.json"

cours_now = {"XLK": 182.5}  # cours du moment (dernier cours disponible)
mon_portefeuille = {}       # lignes détenues, ex. {"XLK": [{"date": "2026-10-20", "prix": 180.1}]}

reglages = dict(train_start="2005-01-01", train_end="2016-12-31", val_end="2019-12-31",
                horizon=HORIZON, target="direction", signaux=True,
                indicator_field_level=0,
                regime_col="score_niveau",
                # score_min=25,
                )
\end{code}

\subsection{Données et indicateurs}
Les prix sont téléchargés avec \f{yfinance} (\f{auto\_adjust=True}) et stockés dans un
dictionnaire \f{prices = \{titre: DataFrame Open/High/Low/Close/Volume\}}. Les indicateurs
sont calculés par \f{construire\_indicateurs}, qui appelle le code d'indicateurs du projet
(\f{construire\_panels}) et renomme les colonnes selon les noms attendus par
\f{signaux.py}. L'ATR est extrait à part pour le stop. Le score n'est pas mis parmi les
indicateurs : il passe par l'argument \f{macro} de \f{prepare\_dataset}.

\begin{code}[indicateurs]{script\_xlk.py}
def construire_indicateurs(prices_):
    pan = construire_panels(prices_, Indicateurs, ParamsPF(min_titres=1),
                            generateur=features_completes)
    ind = {"rsi": pan["rsi"], "bb_pct": pan["bollinger_pct_b"], "rev_5": pan["rev_5"],
           "stoch": pan["stoch"], "er": pan["er"], "hist": pan["hist"]}
    return pd.concat(ind, axis=1, names=["Indicator", "Ticker"]), pan["atr"]
\end{code}

\noindent Pour ajouter un indicateur : le calculer dans le code d'indicateurs, l'ajouter
au dictionnaire \f{ind} sous un nom, puis ajouter les règles correspondantes dans
\f{SIGNAUX} avec ce même nom dans \f{col}.

\subsection{Bloc 1 : entraînement, seuils, test (une fois)}
Exécuté avec \f{ENTRAINER = True}, sur les clôtures :
\begin{enumerate}[nosep]
  \item \f{prepare\_dataset} : variables, cible, découpage (section~\ref{sec:decoupage}) ;
  \item \f{make\_loaders} : fenêtres et batchs ;
  \item \f{train} : entraînement, meilleurs poids dans \f{modele\_proba.pt} ;
  \item \f{rapport} : diagnostic sur la validation (dossier \f{resultats\_proba}) ;
  \item \f{choisir\_seuils\_portefeuille} : seuils sur la validation, enregistrés dans \f{seuils.json} ;
  \item \f{backtest\_portefeuille} et \f{rapport\_portefeuille} : test, regardé une seule fois
  (dossier \f{backtest\_etf\_tech}).
\end{enumerate}

\paragraph{Lire le rapport de validation.} Le dossier \f{resultats\_proba} contient un
fichier Excel et des graphiques : AUC (0,5 = hasard ; 0,52 à 0,55 est réaliste pour ce
type de problème), gain de log-loss par rapport à une prévision constante, courbe de
calibration (fréquence de hausse observée par tranche de probabilité), importance moyenne
des variables (poids du VSN) et lecture par régime de score. Un gain de log-loss positif
ne garantit pas une stratégie rentable : il peut venir des jours où la stratégie
n'achète pas, et la cible mesure le sens de variation, pas son ampleur.

\subsection{Bloc 2 : décisions du jour}
Exécuté à chaque lancement, au moment de la décision. Les données téléchargées s'arrêtent
à la dernière clôture ; \f{figer\_seance} ajoute la séance du jour avec le cours du moment
comme clôture, puis les indicateurs sont recalculés sur ces prix. Le modèle et les seuils
sont relus depuis les fichiers du bloc 1.

\begin{code}[décisions du jour]{script\_xlk.py}
prices_jour = figer_seance(prices, cours_now)          # Close du jour = cours du moment
indicateurs_jour, atr_jour = construire_indicateurs(prices_jour)
data_jour = prepare_dataset(prices_jour, indicateurs_jour, score, **reglages)  # mêmes reglages
model, _ = load_model(FICHIER_MODELE)
s = json.load(open(FICHIER_SEUILS))

dec = decisions_portefeuille_du_jour(model, data_jour, FENETRE, mon_portefeuille,
                                     s["seuil_achat"], s["seuil_sortie"],
                                     prices=prices_jour, **options(atr_jour))
print(dec["vendre"].to_string()); print(dec["garder"].to_string()); print(dec["acheter"].to_string())
\end{code}

\noindent Le résultat contient quatre tableaux : \f{vendre} (lignes à vendre et motif),
\f{garder}, \f{acheter} (nouvelle ligne ou deuxième ligne) et \f{toutes les probabilités}.
La colonne \f{date des données} doit afficher la date du jour pour tous les titres ; un
titre absent de \f{cours\_now} est évalué sur sa dernière clôture. Après avoir passé les
ordres, on met à jour \f{mon\_portefeuille} :
\begin{code}[format du portefeuille]{script\_xlk.py}
mon_portefeuille = {"CAT":   [{"date": "2026-10-20", "prix": 352.1},     # deux lignes
                              {"date": "2026-10-26", "prix": 360.4}],
                    "MC.PA": {"date": "2026-10-15", "prix": 698.5}}      # une ligne
\end{code}

\subsection{Règles de cohérence}
\begin{itemize}[nosep]
  \item \f{reglages}, \f{FENETRE}, les options de la stratégie, le fichier \f{.pt} et
  \f{seuils.json} doivent être identiques entre le bloc 1 et le bloc 2.
  \item Toute modification de \f{reglages}, de \f{signaux.py} ou des indicateurs impose de
  réentraîner (bloc 1). Une modification des seuls paramètres de la stratégie impose de
  refaire le choix des seuils (le modèle peut être rechargé avec \f{load\_model} au lieu
  d'appeler \f{train}).
  \item \f{figer\_seance} met le volume du jour à NaN (volume partiel). Si un indicateur ou
  \f{volume\_rel} en dépend, utiliser \f{volume="garder"} ou \f{variables\_prix=False}.
  \item Après le remplacement d'un module, redémarrer le noyau Python pour éviter de
  mélanger deux versions.
  \item Hors Jupyter, afficher les tableaux avec \f{print(df.to\_string())}.
\end{itemize}

% ===========================================================================
\section{Récapitulatif des paramètres}\label{sec:parametres}
% ===========================================================================
\small
\begin{longtable}{@{}>{\raggedright\arraybackslash}p{3.3cm} >{\raggedright\arraybackslash}p{4.3cm} >{\raggedright\arraybackslash}p{2.1cm} >{\raggedright\arraybackslash}p{5.0cm}@{}}
\toprule
\textbf{Paramètre} & \textbf{Où le modifier} & \textbf{Valeur} & \textbf{Effet} \\
\midrule
\endhead
\multicolumn{4}{@{}l}{\emph{Données et cible}}\\
\f{TICKERS}        & en-tête du script & \f{["XLK"]} & univers de titres \\
\f{HORIZON}        & en-tête du script & 22 & $H$ ; réentraîner \\
\f{FENETRE}        & en-tête du script & 20 & $T$ ; réentraîner \\
\f{train\_start}, \f{train\_end}, \f{val\_end} & \f{reglages} & 2005 / 2016 / 2019 & découpage temporel ; réentraîner \\
\f{target}         & \f{reglages} & \f{"direction"} & cible binaire ($y$) \\
\midrule
\multicolumn{4}{@{}l}{\emph{Signaux et variables}}\\
seuils, sens, colonnes & \f{SIGNAUX} (\f{signaux.py}) & --- & règles de franchissement ; réentraîner \\
\f{demi\_vie}      & \f{reglages} & 2 & vitesse de décroissance de la fraîcheur \\
\f{max\_jours}     & \f{reglages} & 10 & fraîcheur remise à 0 après ce délai \\
\f{garder\_niveaux} & \f{reglages} & \f{False} & ajoute les niveaux bruts des indicateurs \\
\f{variables\_prix} & \f{reglages} & \f{True} & rendements, volatilité, volume relatif \\
\midrule
\multicolumn{4}{@{}l}{\emph{Score macroéconomique}}\\
\f{regime\_col}    & \f{reglages} & \f{"score\_niveau"} & score utilisé comme contexte \\
\f{seuils\_regime} & \f{reglages} & (25, 75) & indicatrices de régime bas / élevé \\
\f{score\_min}, \f{score\_max} & \f{reglages} & désactivé & pas d'apprentissage ni d'achat hors domaine \\
\f{filtre\_sur}    & \f{reglages} & toutes périodes & périodes concernées par le filtre \\
\f{publication\_lag\_days} & \f{load\_macro(...)} & 30 & délai avant utilisation du score \\
\midrule
\multicolumn{4}{@{}l}{\emph{Réseau et entraînement (\f{TrainConfig})}}\\
\f{d\_model}       & \f{TrainConfig} & 16 & taille des représentations internes \\
\f{n\_layers}      & \f{TrainConfig} & 1 & couches de LSTM \\
\f{dropout}        & \f{TrainConfig} & 0,2 & régularisation \\
\f{lr}, \f{weight\_decay} & \f{TrainConfig} & 1e-3, 1e-5 & pas d'apprentissage, pénalité L2 \\
\f{patience}, \f{max\_epochs} & \f{TrainConfig} & 8, 100 & arrêt anticipé \\
\f{batch\_size}    & \f{make\_loaders} & 256 & fenêtres par ajustement \\
\midrule
\multicolumn{4}{@{}l}{\emph{Seuils}}\\
\f{grille\_achat}, \f{grille\_sortie} & \f{choisir\_seuils\_portefeuille} & quantiles & seuils candidats \\
\f{min\_trades}    & \f{choisir\_seuils\_portefeuille} & 10 & trades minimum pour qu'un couple compte \\
\midrule
\multicolumn{4}{@{}l}{\emph{Stratégie (fonction \f{options} du script)}}\\
\f{max\_positions} & \f{options} & $\min(10, 2n)$ & lignes simultanées ; taille $=1/$\f{max\_positions} \\
\f{max\_par\_titre} & \f{options} & 2 & lignes par titre \\
\f{ecart\_min}     & \f{options} & 3 & séances entre deux achats d'un même titre \\
\f{recroisement}   & \f{options} & \f{False} & 2\textsuperscript{e} ligne seulement après un repli sous le seuil \\
\f{max\_entrees\_jour} & \f{options} & 2 & achats par jour \\
\f{stop\_atr}      & \f{options} & 2 & stop en ATR ; \f{None} = désactivé \\
\f{duree\_max}     & \f{options} & \f{None} & sortie au temps ; \f{None} = désactivée \\
\f{couts\_bps}     & \f{options} & 5 & frais par transaction \\
\f{entree}         & \f{options} & \f{"cloture"} & exécution au cours du calcul ou à l'ouverture suivante \\
\f{taille}         & \f{options} (à ajouter) & \f{"courant"} & base de calcul de la taille des lignes \\
\midrule
\multicolumn{4}{@{}l}{\emph{Usage quotidien}}\\
\f{ENTRAINER}      & en-tête du script & \f{True} & \f{False} : décisions du jour seulement \\
\f{cours\_now}     & en-tête du script & --- & cours du moment par titre \\
\f{mon\_portefeuille} & en-tête du script & --- & lignes détenues (date, prix) \\
\bottomrule
\end{longtable}
\normalsize

\end{document}
